% GENERATED FILE. Edit canonical lessons, then run npm run generate:books.
% canonical-source-hash: fnv1a64:ac9d0a9d834897b0
% canonical-lessons: ZH-C224-kaihua, ZH-W224-hua-flower, ZH-C224-huadian, ZH-C224-huayuan, ZH-C224-hudie, ZH-C224-mifeng

\chapter{Spring Flowers}
\label{ch:zh-spring-flowers}
\hlchaptermodality{224}

\begin{chapteropening}
\textbf{By the end of this chapter:} \emph{I can say flowers are coming out, and talk about a florist's, a garden, butterflies and bees.}

Complete the last lesson of chapter 224: I can say flowers are coming out, and talk about a florist's, a garden, butterflies and bees.
\end{chapteropening}

\section[\texorpdfstring{\zh{开花} (kāihuā)}{开花 (kāihuā)}]{\zh{开花} (kāihuā) — to bloom; to flower}

\label{lesson:ZH-C224-kaihua}



\begin{quote}
Before the new one: say the Mandarin for a monkey, then the Mandarin for a rabbit.
\end{quote}

\subsection*{You'll want to know: \zh{开花}}
\textbf{\zh{开花}} — \emph{kāihuā} — \textquotedblleft{}to bloom; to flower\textquotedblright{}.

\begin{quote}
\textbf{\zh{这些花三月开花}} \emph{zhèxiē huā sānyuè kāihuā} — These flowers come out in March.
\end{quote}

The new character, in pieces: \zh{花} — seven strokes: the grass piece of \zh{菜} on top, then the person piece, then two strokes like the right half of \zh{比}. The next lesson puts it in your hand.

\begin{cousinweb}
\textbf{\zh{开花}} (\emph{kāihuā}) is built from \zh{开} 'to open' and \zh{花} 'flower', both inherited from Middle Chinese, \emph{khoj} and \emph{xwae} (Guangyun, in Baxter's transcription), which notes the second as the everyday form of the word for blossom.
\end{cousinweb}

\begin{grammarlens}[title={What you've built}]
One more word for describing things.
\end{grammarlens}

\subsection*{Your turn}
\begin{itemize}
\raggedright
  \item \emph{Say it:} \emph{kāihuā}
  \item \emph{Say it:} \emph{kāihuā}, once more
  \item \emph{Say it:} \emph{zhèxiē huā sānyuè kāihuā} — These flowers come out in March.
  \item \emph{Recall:} say the Mandarin for a monkey, then the Mandarin for a rabbit, then say \emph{kāihuā} again
\end{itemize}

\begin{tcolorbox}[breakable,colback=teal!4,colframe=teal!35!black,title={Before you move on}]
What does \zh{开花} mean? (\textquotedblleft{}To bloom; to flower\textquotedblright{}.) Say it once more.
\end{tcolorbox}
\section[\texorpdfstring{\zh{花} (huā)}{花 (huā)}]{\zh{花} — the grass piece over the person piece and two more strokes}

\label{lesson:ZH-W224-hua-flower}



\begin{quote}
The word first, then the character inside it.

\begin{itemize}
\raggedright
  \item \emph{Recall:} say \emph{kāihuā} — \textbf{R1}, one lesson back, before it decays
\end{itemize}

\begin{quote}
\zh{菜} \zh{比}
\end{quote}

\emph{Read it:} each one aloud, with nothing said first — \emph{cài}, \emph{bǐ}
\end{quote}

\subsection*{Script — the grass piece over the person piece and two more strokes}
\hlblockfigure{\detokenize{figures/ZH-W224-hua-flower-filmstrip.pdf}}{How \zh{花} is written, stroke by stroke}

\begin{quote}
\zh{花}
\end{quote}

\textbf{Three pieces: one on top, two side by side under it.} Count the pieces, not the strokes.

\textbf{The grass piece, three strokes,} as in \zh{菜} and \zh{茶}: a bar across, then two short verticals through it.

\textbf{The person piece, \zh{亻}, two strokes,} as in \zh{你}: a stroke falling to the left, then a vertical.

\textbf{The right piece, two strokes,} like the right half of \zh{比}: a short stroke falling to the left, then a vertical that bends right along the foot and flicks up.

3 + 2 + 2 = seven.

\subsection*{Your turn}
\begin{itemize}
\raggedright
  \item \emph{Write it:} \zh{花}, all seven, with the model in front of you
  \item \emph{Cover:} write it again from memory, and count the pieces, not the strokes
\end{itemize}

\begin{tcolorbox}[breakable,colback=teal!4,colframe=teal!35!black,title={Before you move on}]
How many strokes in \zh{花}? (\textbf{Seven}.) What is on top? (\textbf{The grass piece of \zh{菜}}.) Which piece comes next? (\textbf{\zh{亻}, the person piece}.)

Source: \href{https://raw.githubusercontent.com/chanind/hanzi-writer-data/68d10a4b21150cae5e1ebbd223eed289cf32d90c/data/\%E8\%8A\%B1.json}{Hanzi Writer Data \zh{花}, pinned stroke paths and medians}.
\end{tcolorbox}
\section[\texorpdfstring{\zh{花店} (huādiàn)}{花店 (huādiàn)}]{\zh{花店} (huādiàn) — florist's; flower shop}

\label{lesson:ZH-C224-huadian}



\begin{quote}
Before the new one: say the Mandarin for a rabbit, then the Mandarin for to bloom.
\end{quote}

\subsection*{You'll want to know: \zh{花店}}
\textbf{\zh{花店}} — \emph{huādiàn} — \textquotedblleft{}florist's; flower shop\textquotedblright{}.

\begin{quote}
\textbf{\zh{我去花店买花}} \emph{wǒ qù huādiàn mǎi huā} — I'm going to the florist's to buy flowers.
\end{quote}

\begin{cousinweb}
\textbf{\zh{花店}} (\emph{huādiàn}) is built from \zh{花} 'flower' and \zh{店} 'shop', both inherited from Middle Chinese, \emph{xwae} and \emph{temH} (Guangyun, in Baxter's transcription).
\end{cousinweb}

\begin{grammarlens}[title={What you've built}]
One more word for describing things.
\end{grammarlens}

\subsection*{Your turn}
\begin{itemize}
\raggedright
  \item \emph{Say it:} \emph{huādiàn}
  \item \emph{Say it:} \emph{huādiàn}, once more
  \item \emph{Say it:} \emph{wǒ qù huādiàn mǎi huā} — I'm going to the florist's to buy flowers.
  \item \emph{Recall:} say the Mandarin for a rabbit, then the Mandarin for to bloom, then say \emph{huādiàn} again
\end{itemize}

\begin{tcolorbox}[breakable,colback=teal!4,colframe=teal!35!black,title={Before you move on}]
What does \zh{花店} mean? (\textquotedblleft{}Florist's; flower shop\textquotedblright{}.) Say it once more.
\end{tcolorbox}
\section[huāyuán]{huāyuán — garden}

\label{lesson:ZH-C224-huayuan}



\begin{quote}
Before the new one: say the Mandarin for to bloom, then the Mandarin for a florist's.
\end{quote}

\subsection*{You'll want to know: huāyuán}
\emph{huāyuán} — \textquotedblleft{}garden\textquotedblright{}.

\begin{quote}
\emph{wǒ jiā hòubian yǒu yí ge huāyuán} — There's a garden behind our house.
\end{quote}

Its characters come later; for now it is a word for the ear.

\begin{cousinweb}
\emph{huāyuán} is built from \zh{花} 'flower' and \emph{yuán} 'a plot', both inherited from Middle Chinese, \emph{xwae} and \emph{hjwon} (Guangyun, in Baxter's transcription), which glosses the second 'a garden plot'.
\end{cousinweb}

\begin{grammarlens}[title={What you've built}]
One more word for describing things.
\end{grammarlens}

\subsection*{Your turn}
\begin{itemize}
\raggedright
  \item \emph{Say it:} \emph{huāyuán}
  \item \emph{Say it:} \emph{huāyuán}, once more
  \item \emph{Say it:} \emph{wǒ jiā hòubian yǒu yí ge huāyuán} — There's a garden behind our house.
  \item \emph{Recall:} say the Mandarin for to bloom, then the Mandarin for a florist's, then say \emph{huāyuán} again
\end{itemize}

\begin{tcolorbox}[breakable,colback=teal!4,colframe=teal!35!black,title={Before you move on}]
What does huāyuán mean? (\textquotedblleft{}Garden\textquotedblright{}.) Say it once more.
\end{tcolorbox}
\section[húdié]{húdié — butterfly}

\label{lesson:ZH-C224-hudie}



\begin{quote}
Before the new one: say the Mandarin for a florist's, then the Mandarin for a garden.
\end{quote}

\subsection*{You'll want to know: húdié}
\emph{húdié} — \textquotedblleft{}butterfly\textquotedblright{}.

\begin{quote}
\emph{huāyuán lǐ yǒu liǎng zhī húdié} — There are two butterflies in the garden.
\end{quote}

Its characters come later; for now it is a word for the ear.

\begin{cousinweb}
\emph{húdié} is built from two characters, and the dictionary glosses the second, \emph{dep} there (Guangyun, in Baxter's transcription), as a butterfly. No dictionary we use records an older form of the first, so its history is not settled.
\end{cousinweb}

\begin{grammarlens}[title={What you've built}]
One more word for describing things.
\end{grammarlens}

\subsection*{Your turn}
\begin{itemize}
\raggedright
  \item \emph{Say it:} \emph{húdié}
  \item \emph{Say it:} \emph{húdié}, once more
  \item \emph{Say it:} \emph{huāyuán lǐ yǒu liǎng zhī húdié} — There are two butterflies in the garden.
  \item \emph{Recall:} say the Mandarin for a florist's, then the Mandarin for a garden, then say \emph{húdié} again
\end{itemize}

\begin{tcolorbox}[breakable,colback=teal!4,colframe=teal!35!black,title={Before you move on}]
What does húdié mean? (\textquotedblleft{}Butterfly\textquotedblright{}.) Say it once more.
\end{tcolorbox}
\section[mìfēng]{mìfēng — bee; honeybee}

\label{lesson:ZH-C224-mifeng}



\begin{quote}
Before the new one: say the Mandarin for a garden, then the Mandarin for a butterfly.
\end{quote}

\subsection*{You'll want to know: mìfēng}
\emph{mìfēng} — \textquotedblleft{}bee; honeybee\textquotedblright{}.

\begin{quote}
\emph{huāyuán lǐ de huā kāi le mìfēng yě lái le} — The flowers in the garden are out, and the bees have come too.
\end{quote}

Its characters come later; for now it is a word for the ear.

\begin{cousinweb}
\emph{mìfēng} is built from \emph{mì} 'honey' and \emph{fēng} 'bee', both inherited from Middle Chinese, \emph{mjit} and \emph{phjowng} (Guangyun, in Baxter's transcription), which glosses the first as the food bees make and explains the second, after an older book, as a flying insect that stings.
\end{cousinweb}

\begin{grammarlens}[title={What you've built}]
One more word for describing things.
\end{grammarlens}

\subsection*{Your turn}
\begin{itemize}
\raggedright
  \item \emph{Say it:} \emph{mìfēng}
  \item \emph{Say it:} \emph{mìfēng}, once more
  \item \emph{Say it:} \emph{huāyuán lǐ de huā kāi le mìfēng yě lái le} — The flowers in the garden are out, and the bees have come too.
  \item \emph{Recall:} say the Mandarin for a garden, then the Mandarin for a butterfly, then say \emph{mìfēng} again
\end{itemize}

\begin{tcolorbox}[breakable,colback=teal!4,colframe=teal!35!black,title={Before you move on}]
What does mìfēng mean? (\textquotedblleft{}Bee; honeybee\textquotedblright{}.) Say it once more.
\end{tcolorbox}
